\section{Discussion}
\label{sec:discussion}

\paragraph{Where residual mass remains.}
Across the 100\,000-cycle procedural bench and lit-combo family stress tests, wrap error is not uniform.
The same generator families repeatedly dominate the residual budget: \texttt{nonlinear}, \texttt{combo}, \texttt{triple\_mix}, and \texttt{extreme\_overlay} show the weakest denoise / residual scores, while \texttt{harmonic\_fft} and \texttt{am\_fm} are comparatively easy.
\texttt{open\_wrap\_bias} can look strong on wrap-energy proxies yet still lag on prolonged residual---a reminder that cliff reduction and ideal-matched tiling are related but not identical objectives.

\paragraph{Implication for this paper.}
Mean $R$ (and mean DualCosine gaps) therefore compress a structured difficulty map.
Reporting only the overnight mean on synthetic sine cliffs understates how much of the remaining $\sim$1\% sits in a few hard families.
That structure is a first-class empirical claim of the present work, independent of whether the GPU hybrid later nudges mean $R$ from $\approx 0.99$ toward $0.999$.

\paragraph{Hybrid search vs.\ plateau.}
Interim overnight traces show early gains then long plateaus near $R\approx 0.99$.
Plateau adaptation (deeper cells, denser MoE graphs, higher GA/PBT weight) is the operational response inside one campaign; it does not by itself rebalance family-conditional risk.
A natural next paper is family-conditional and cliff-conditional meta-learning: train or select denoisers per family, or condition the outer loop on measured wrap statistics, and publish per-family ROC-style residual curves rather than a single scalar champion.

\paragraph{Hybrid branches (interim).}
Final branch win-rates (GA / PPO / PBT / NAS / lit-combo) remain deferred until the declared GPU budget completes; anticipated comparison is against the v3 finding that light evolutionary explore beat nested loss-opt elites under a smaller residual budget.
